% GENERATED FILE. Edit canonical lessons, then run npm run generate:books.
% canonical-source-hash: fnv1a64:53ad0939a71c2aa1
% canonical-lessons: ML-C362-tipitittam, ML-C362-agnisamanasena, ML-C362-campal, ML-C362-tipitikkuka, ML-C362-kattinasikkuka

\chapter{A Fire in the Market}
\label{ch:ml-a-fire-in-the-market}
\hlchaptermodality{362}

\begin{chapteropening}
\textbf{By the end of this chapter:} \emph{I can say a fire, the fire brigade, ashes, to catch fire and to burn down.}

Complete the last lesson of chapter 362: I can say a fire, the fire brigade, ashes, to catch fire and to burn down.
\end{chapteropening}

\section[\texorpdfstring{\ml{തീപിടിത്തം} (tīpiṭittaṁ)}{തീപിടിത്തം (tīpiṭittaṁ)}]{\ml{തീപിടിത്തം} (tīpiṭittaṁ) — a fire; a blaze}

\label{lesson:ML-C362-tipitittam}



\begin{quote}
Before the new one: say the Malayalam for a beginning, then the Malayalam for unexpectedly.
\end{quote}

\subsection*{You'll want to know: \ml{തീപിടിത്തം}}
\textbf{\ml{തീപിടിത്തം}} — \emph{tīpiṭittaṁ} — \textquotedblleft{}a fire; a blaze\textquotedblright{}.

\begin{quote}
\ml{ഇന്നലെ} \ml{രാത്രി} \ml{ചന്തയിൽ} \ml{ഒരു} \ml{വലിയ} \ml{തീപിടിത്തം} \ml{ഉണ്ടായി}.
\end{quote}

\begin{quote}
\emph{innale rātri cantayil oru valiya tīpiṭittaṁ uṇṭāyi.}
\end{quote}

\begin{quote}
Last night there was a big fire at the market.
\end{quote}

Literally a catching of fire: \ml{തീ}, fire, + \ml{പിടിത്തം}, a catching. A news word for a fire that breaks out.

\begin{cousinweb}
\textbf{\ml{തീപിടിത്തം}} (\emph{tīpiṭittaṁ}) is a native Dravidian compound of \ml{തീ} (\emph{tī}), \textquotedblleft{}fire\textquotedblright{} (DEDR 3266), and \ml{പിടിത്തം} (\emph{piṭittaṁ}), \textquotedblleft{}a catching, a grip\textquotedblright{}, from \ml{പിടിക്കുക}, \textquotedblleft{}to catch\textquotedblright{} (DEDR 4148). A building catches fire.
\end{cousinweb}

\begin{grammarlens}[title={What you've built}]
One more word for telling what happened.
\end{grammarlens}

\subsection*{Your turn}
\begin{itemize}
\raggedright
  \item \emph{Say it:} \emph{tīpiṭittaṁ}
  \item \emph{Say it:} \emph{tīpiṭittaṁ}, once more
  \item \emph{Say it:} \emph{innale rātri cantayil oru valiya tīpiṭittaṁ uṇṭāyi.} — Last night there was a big fire at the market.
  \item \emph{Recall:} say the Malayalam for a beginning, then the Malayalam for unexpectedly, then say \emph{tīpiṭittaṁ} again
\end{itemize}

\begin{tcolorbox}[breakable,colback=teal!4,colframe=teal!35!black,title={Before you move on}]
What does \ml{തീപിടിത്തം} mean? (\textquotedblleft{}A fire; a blaze\textquotedblright{}.) Say it once more.
\end{tcolorbox}
\section[\texorpdfstring{\ml{അഗ്നിശമനസേന} (agniśamanasēna)}{അഗ്നിശമനസേന (agniśamanasēna)}]{\ml{അഗ്നിശമനസേന} (agniśamanasēna) — the fire brigade}

\label{lesson:ML-C362-agnisamanasena}



\begin{quote}
Before the new one: say the Malayalam for unexpectedly, then the Malayalam for a fire.
\end{quote}

\subsection*{You'll want to know: \ml{അഗ്നിശമനസേന}}
\textbf{\ml{അഗ്നിശമനസേന}} — \emph{agniśamanasēna} — \textquotedblleft{}the fire brigade\textquotedblright{}.

\begin{quote}
\ml{അഗ്നിശമനസേന} \ml{അര} \ml{മണിക്കൂറിനുള്ളിൽ} \ml{എത്തി}.
\end{quote}

\begin{quote}
\emph{agniśamanasēna ara maṇikkūṟinuḷḷil etti.}
\end{quote}

\begin{quote}
The fire brigade arrived within half an hour.
\end{quote}

Three Sanskrit words in a row: fire, quenching, force.

\begin{cousinweb}
\textbf{\ml{അഗ്നിശമനസേന}} (\emph{agniśamanasēna}) is built from three Sanskrit words: \emph{agni}, \textquotedblleft{}fire\textquotedblright{} (CDIAL 55), \emph{śamana}, \textquotedblleft{}calming, extinguishing\textquotedblright{}, and \emph{senā}, \textquotedblleft{}an army\textquotedblright{} (Monier-Williams). It is the force that puts fires out.
\end{cousinweb}

\begin{grammarlens}[title={What you've built}]
One more word for telling what happened.
\end{grammarlens}

\subsection*{Your turn}
\begin{itemize}
\raggedright
  \item \emph{Say it:} \emph{agniśamanasēna}
  \item \emph{Say it:} \emph{agniśamanasēna}, once more
  \item \emph{Say it:} \emph{agniśamanasēna ara maṇikkūṟinuḷḷil etti.} — The fire brigade arrived within half an hour.
  \item \emph{Recall:} say the Malayalam for unexpectedly, then the Malayalam for a fire, then say \emph{agniśamanasēna} again
\end{itemize}

\begin{tcolorbox}[breakable,colback=teal!4,colframe=teal!35!black,title={Before you move on}]
What does \ml{അഗ്നിശമനസേന} mean? (\textquotedblleft{}The fire brigade\textquotedblright{}.) Say it once more.
\end{tcolorbox}
\section[\texorpdfstring{\ml{ചാമ്പൽ} (cāmpal)}{ചാമ്പൽ (cāmpal)}]{\ml{ചാമ്പൽ} (cāmpal) — ashes; cinders}

\label{lesson:ML-C362-campal}



\begin{quote}
Before the new one: say the Malayalam for a fire, then the Malayalam for the fire brigade.
\end{quote}

\subsection*{You'll want to know: \ml{ചാമ്പൽ}}
\textbf{\ml{ചാമ്പൽ}} — \emph{cāmpal} — \textquotedblleft{}ashes; cinders\textquotedblright{}.

\begin{quote}
\ml{തീപിടിത്തത്തിന്} \ml{ശേഷം} \ml{കടകൾ} \ml{ചാമ്പലായി}.
\end{quote}

\begin{quote}
\emph{tīpiṭittattinŭ śēṣaṁ kaṭakaḷ cāmpalāyi.}
\end{quote}

\begin{quote}
After the fire the shops were reduced to ashes.
\end{quote}

\ml{ചാമ്പലാകുക}, to become ashes, to burn to the ground. \ml{ചാരം} is the everyday ash of a stove.

\begin{cousinweb}
\textbf{\ml{ചാമ്പൽ}} (\emph{cāmpal}) is native Dravidian: DEDR 2453 lists Malayalam \emph{cāmpal}, \textquotedblleft{}closing; ashes\textquotedblright{}, under \ml{ചാമ്പുക}, \textquotedblleft{}to shut, to wither\textquotedblright{}: what is left when a fire dies down.
\end{cousinweb}

\begin{grammarlens}[title={What you've built}]
One more word for telling what happened.
\end{grammarlens}

\subsection*{Your turn}
\begin{itemize}
\raggedright
  \item \emph{Say it:} \emph{cāmpal}
  \item \emph{Say it:} \emph{cāmpal}, once more
  \item \emph{Say it:} \emph{tīpiṭittattinŭ śēṣaṁ kaṭakaḷ cāmpalāyi.} — After the fire the shops were reduced to ashes.
  \item \emph{Recall:} say the Malayalam for a fire, then the Malayalam for the fire brigade, then say \emph{cāmpal} again
\end{itemize}

\begin{tcolorbox}[breakable,colback=teal!4,colframe=teal!35!black,title={Before you move on}]
What does \ml{ചാമ്പൽ} mean? (\textquotedblleft{}Ashes; cinders\textquotedblright{}.) Say it once more.
\end{tcolorbox}
\section[\texorpdfstring{\ml{തീപിടിക്കുക} (tīpiṭikkuka)}{തീപിടിക്കുക (tīpiṭikkuka)}]{\ml{തീപിടിക്കുക} (tīpiṭikkuka) — to catch fire}

\label{lesson:ML-C362-tipitikkuka}



\begin{quote}
Before the new one: say the Malayalam for the fire brigade, then the Malayalam for ashes.
\end{quote}

\subsection*{You'll want to know: \ml{തീപിടിക്കുക}}
\textbf{\ml{തീപിടിക്കുക}} — \emph{tīpiṭikkuka} — \textquotedblleft{}to catch fire\textquotedblright{}.

\begin{quote}
\ml{അടുക്കളയിലെ} \ml{ഗ്യാസ്} \ml{കാരണം} \ml{വീടിന്} \ml{തീപിടിച്ചു}.
\end{quote}

\begin{quote}
\emph{aṭukkaḷayile gyāsŭ kāraṇaṁ vīṭinŭ tīpiṭiccu.}
\end{quote}

\begin{quote}
The house caught fire because of the gas in the kitchen.
\end{quote}

Past: \ml{തീപിടിച്ചു}, \emph{tīpiṭiccu}.

What catches fire takes -\ml{ക്ക്} or -\ml{ന്}: \ml{കടയ്ക്ക്} \ml{തീപിടിച്ചു}, the shop caught fire.

\begin{cousinweb}
\textbf{\ml{തീപിടിക്കുക}} (\emph{tīpiṭikkuka}) is a native Dravidian compound of \ml{തീ} (\emph{tī}), \textquotedblleft{}fire\textquotedblright{} (DEDR 3266), and \ml{പിടിക്കുക} (\emph{piṭikkuka}), \textquotedblleft{}to catch, to seize\textquotedblright{} (DEDR 4148). Its noun is \ml{തീപിടിത്തം}.
\end{cousinweb}

\begin{grammarlens}[title={What you've built}]
One more word for telling what happened.
\end{grammarlens}

\subsection*{Your turn}
\begin{itemize}
\raggedright
  \item \emph{Say it:} \emph{tīpiṭikkuka}
  \item \emph{Say it:} \emph{tīpiṭikkuka}, once more
  \item \emph{Say it:} \emph{aṭukkaḷayile gyāsŭ kāraṇaṁ vīṭinŭ tīpiṭiccu.} — The house caught fire because of the gas in the kitchen.
  \item \emph{Recall:} say the Malayalam for the fire brigade, then the Malayalam for ashes, then say \emph{tīpiṭikkuka} again
\end{itemize}

\begin{tcolorbox}[breakable,colback=teal!4,colframe=teal!35!black,title={Before you move on}]
What does \ml{തീപിടിക്കുക} mean? (\textquotedblleft{}To catch fire\textquotedblright{}.) Say it once more.
\end{tcolorbox}
\section[\texorpdfstring{\ml{കത്തിനശിക്കുക} (kattinaśikkuka)}{കത്തിനശിക്കുക (kattinaśikkuka)}]{\ml{കത്തിനശിക്കുക} (kattinaśikkuka) — to burn down}

\label{lesson:ML-C362-kattinasikkuka}



\begin{quote}
Before the new one: say the Malayalam for ashes, then the Malayalam for to catch fire.
\end{quote}

\subsection*{You'll want to know: \ml{കത്തിനശിക്കുക}}
\textbf{\ml{കത്തിനശിക്കുക}} — \emph{kattinaśikkuka} — \textquotedblleft{}to burn down\textquotedblright{}.

\begin{quote}
\ml{ആ} \ml{പഴയ} \ml{കെട്ടിടം} \ml{മുഴുവൻ} \ml{കത്തിനശിച്ചു}.
\end{quote}

\begin{quote}
\emph{ā paḻaya keṭṭiṭaṁ muḻuvan kattinaśiccu.}
\end{quote}

\begin{quote}
That whole old building burned down.
\end{quote}

Past: \ml{കത്തിനശിച്ചു}, \emph{kattinaśiccu}.

Burned and ruined: \ml{കത്തി}, burning, + a Sanskrit verb of ruin. A news word.

\begin{cousinweb}
\textbf{\ml{കത്തിനശിക്കുക}} (\emph{kattinaśikkuka}) joins the native verb \ml{കത്തുക} (\emph{kattuka}), \textquotedblleft{}to burn\textquotedblright{} (DEDR 1207), and \ml{നശിക്കുക}, \textquotedblleft{}to perish\textquotedblright{}, built from Sanskrit \emph{nāśa}, \textquotedblleft{}destruction, ruin\textquotedblright{} (CDIAL 7084): to burn until nothing is left.
\end{cousinweb}

\begin{grammarlens}[title={What you've built}]
One more word for telling what happened.
\end{grammarlens}

\subsection*{Your turn}
\begin{itemize}
\raggedright
  \item \emph{Say it:} \emph{kattinaśikkuka}
  \item \emph{Say it:} \emph{kattinaśikkuka}, once more
  \item \emph{Say it:} \emph{ā paḻaya keṭṭiṭaṁ muḻuvan kattinaśiccu.} — That whole old building burned down.
  \item \emph{Recall:} say the Malayalam for ashes, then the Malayalam for to catch fire, then say \emph{kattinaśikkuka} again
\end{itemize}

\begin{tcolorbox}[breakable,colback=teal!4,colframe=teal!35!black,title={Before you move on}]
What does \ml{കത്തിനശിക്കുക} mean? (\textquotedblleft{}To burn down\textquotedblright{}.) Say it once more.
\end{tcolorbox}
